% ===================================================================
% SKILL: KAGGLE ÖDEV VE UYGULAMA SAYFASI
% ===================================================================
% TANIM:
% Bu şablon, öğrencilere haftalık Kaggle ödevlerini ve uygulama 
% ortamını tanıtmak için kullanılan standart "Haftalık Ödev ve Uygulama" 
% frame'ini oluşturur.
% 
% ÖZELLİKLER:
% 1. Adım adım talimat listesi (Hesap Açma -> Erişim -> Kopyalama)
% 2. Tıklanabilir ve renklendirilmiş (hrefcol) link yapısı
% 3. Görsel vurgu için ortalanmış buton benzeri link alanı
% ===================================================================

% \section{Haftalık Ödev} % Ana akışta uygun yere eklenebilir

\begin{frame}{Haftalık Ödev ve Uygulama}
    \justifying{}
    Bu haftaki uygulama çalışmaları ve ödevleriniz için aşağıdaki adımları takip ediniz:

    \vspace{1em}
    \begin{enumerate}
        % ADIM 1: HESAP
        \item \textbf{Kaggle Hesabı:} \href{https://www.kaggle.com}{Kaggle} platformunda hesap oluşturunuz (veya giriş yapınız).

        % ADIM 2: DERS BAĞLANTISI (Buradaki URL her hafta değişecek)
        \item \textbf{Derse Erişim:} Aşağıdaki bağlantıyı kullanarak ders not defterine (Notebook) gidiniz:
              \vspace{0.5em}
              \begin{center}
                  % NOT: Aşağıdaki linki haftaya göre güncellemeyi unutmayın!
                  \hrefcol{https://www.kaggle.com/code/ufukasil/mat204-olasilik-ve-stat-st-k-1-hafta}{Kaggle Uygulama Sayfası (Ders Ödevi)}
              \end{center}
              \vspace{0.5em}

        % ADIM 3: ÇALIŞMA ORTAMI (Copy & Edit Vurgusu)
        \item \textbf{Çalışma Ortamı:} Sayfanın sağ üst köşesindeki ``Copy \& Edit'' butonuna tıklayarak kendi kopyanızı oluşturunuz.

        % ADIM 4: TAMAMLAMA
        \item \textbf{Tamamlama:} Defter içerisindeki adımları takip ederek verilen görevleri yerine getiriniz.
    \end{enumerate}
\end{frame}
